\documentclass[11pt,a4paper]{article}
\usepackage[T1]{fontenc}
\usepackage{lmodern}
\usepackage[margin=25mm]{geometry}
\usepackage{amsmath,amssymb,amsthm,mathtools}
\usepackage{microtype,booktabs,array,longtable,xcolor}
\usepackage{listings}
\definecolor{navy}{RGB}{27,54,86}
\definecolor{pale}{RGB}{246,248,250}
\usepackage[colorlinks=true,linkcolor=navy,citecolor=navy,urlcolor=navy]{hyperref}
\usepackage{fancyhdr}
\pagestyle{fancy}\fancyhf{}
\fancyhead[L]{\small Explicit counterexamples to Simon's Le Roy conjecture}
\fancyhead[R]{\small Research manuscript}
\fancyfoot[C]{\thepage}
\setlength{\headheight}{14pt}
\setlength{\emergencystretch}{2em}
\numberwithin{equation}{section}
\newtheorem{theorem}{Theorem}[section]
\newtheorem{proposition}[theorem]{Proposition}
\newtheorem{lemma}[theorem]{Lemma}
\newtheorem{corollary}[theorem]{Corollary}
\theoremstyle{definition}\newtheorem{definition}[theorem]{Definition}
\theoremstyle{remark}\newtheorem{remark}[theorem]{Remark}
\newcommand{\R}{\mathbb R}
\newcommand{\C}{\mathbb C}
\newcommand{\E}{\mathbb E}
\newcommand{\ii}{\mathrm i}
\newcommand{\dd}{\,\mathrm d}
\DeclareMathOperator{\Rea}{Re}
\DeclareMathOperator{\Beta}{Beta}
\newcommand{\bL}{b_{\mathrm L}}
\newcommand{\bU}{b_{\mathrm U}}
\lstset{language=Python,basicstyle=\ttfamily\scriptsize,
 breaklines=true,columns=fullflexible,keepspaces=true,
 showstringspaces=false,frame=single,rulecolor=\color{black!15},
 backgroundcolor=\color{pale},aboveskip=9pt,belowskip=9pt,
 xleftmargin=3pt,xrightmargin=3pt}
\title{Explicit rational counterexamples to\\Simon’s Le Roy conjecture\\[6pt]
\large A beta-product density and a certified complete-monotonicity boundary}
\author{Keio Research Notes\thanks{AI-assisted mathematical research manuscript.
Analytical arguments and executable certificates were checked internally;
this is not external peer review or a claim of historical priority.}}
\date{8 November 2026}
\hypersetup{pdftitle={Explicit rational counterexamples to Simon's Le Roy conjecture},
pdfauthor={Keio Research Notes},pdfsubject={A beta-product proof and certified positivity boundary}}
\begin{filecontents*}[overwrite]{leroy_explicit_certificate.py}
"""Exact Bernstein certificate for explicit Le Roy counterexamples.

Only three reciprocal-gamma evaluations per parameter enclosure use Arb.
Everything after the checked rational enclosures is exact fmpq arithmetic.
Requires python-flint. No floating-point value is used to decide a sign.
"""
import json
import math
from pathlib import Path
import flint
from flint import arb, fmpq as Q, ctx

ctx.prec = 384
N = 90
DIGITS = 80
DEN = 10 ** DIGITS
BL = Q(428292245841, 10**12)
BU = Q(428292245842, 10**12)
records = []


def rational_enclosure(v):
    a = (v.lower() * DEN).floor().unique_fmpz()
    b = (v.upper() * DEN).floor().unique_fmpz()
    assert a is not None and b is not None
    lo, hi = Q(int(a)-1, DEN), Q(int(b)+2, DEN)
    assert arb(lo) < v and v < arb(hi)
    return lo, hi


def coefficients(bl, bu):
    bb = arb(bl).union(arb(bu))
    lo, hi = [], []
    seeds = []
    for r in range(3):
        v = (bb - arb(Q(r+1, 3))).rgamma()**2 / math.factorial(r)
        l, u = rational_enclosure(v)
        assert l > 0
        lo.append(l)
        hi.append(u)
        seeds.append([str(l), str(u)])
    for n in range(N):
        den = 9*(n+1)*(n+2)*(n+3)
        lo.append(lo[n]*(n+4-3*bu)**2/den)
        hi.append(hi[n]*(n+4-3*bl)**2/den)
    return lo, hi, seeds


def falling(n, d):
    return math.factorial(n)//math.factorial(n-d)


def tail_bound(hi, d):
    # For n>=90 and 1/3<beta<1/2, the three-step term ratio,
    # even after d<=2 derivatives, is at most the expression below.
    assert 0 <= d <= 2
    r = Q(64*(N+3)**2, 9*(N+1)*N*(N-1))
    assert r < Q(1, 10)
    return sum(hi[n]*falling(n, d)*4**(n-d)
               for n in range(N, N+3))*Q(10, 9)


def lower_polynomial(lo, hi, d=0):
    p = [((-1)**n)*(lo[n] if n % 2 == 0 else hi[n])*falling(n, d)
         for n in range(d, N)]
    e = tail_bound(hi, d)
    p[0] -= e
    return p, e


def bernstein_on_interval(p, a, b):
    # p(a+(b-a)t)=sum a_k t^k=sum b_j binom(m,j)t^j(1-t)^(m-j).
    m = len(p)-1
    aa = [a**i for i in range(m+1)]
    width = b-a
    power = [sum(p[k]*math.comb(k,j)*aa[k-j]
                 for k in range(j,m+1))*width**j for j in range(m+1)]
    return [sum(power[k]*Q(math.comb(j,k),math.comb(m,k))
                for k in range(j+1)) for j in range(m+1)]


def split_bernstein(b):
    row = b[:]
    left, right = [row[0]], [row[-1]]
    while len(row)>1:
        row = [(row[i]+row[i+1])/2 for i in range(len(row)-1)]
        left.append(row[0])
        right.append(row[-1])
    return left, list(reversed(right))


def certify_positive(p, a, b, target, name):
    initial = bernstein_on_interval(p, a, b)
    stack = [(a,b,initial,0)]
    leaves = []
    minimum = None
    while stack:
        a0,b0,bb,depth = stack.pop()
        lb = min(bb)
        if lb > target:
            leaves.append([str(a0),str(b0),depth])
            minimum = lb if minimum is None else min(minimum,lb)
            continue
        assert depth < 32, (name, str(a0), str(b0), str(lb))
        left,right = split_bernstein(bb)
        mid = (a0+b0)/2
        stack.append((mid,b0,right,depth+1))
        stack.append((a0,mid,left,depth+1))
    # Explicit check that the certified closed intervals cover [a,b].
    assert Q(leaves[0][0]) == a and Q(leaves[-1][1]) == b
    assert all(leaves[i][1] == leaves[i+1][0] for i in range(len(leaves)-1))
    print(name, ':', len(leaves), 'leaves; max depth', max(t[2] for t in leaves),
          '; exact lower bound exceeds', target, flush=True)
    records.append({'name':name,'interval':[str(a),str(b)],'target':str(target),
                    'leaves':leaves,'minimum_coefficient':str(minimum)})


def eval_poly(p, x):
    ans = Q(0)
    for a in reversed(p):
        ans = ans*x+a
    return ans


def main():
    print('python-flint',flint.__version__,'; Arb precision',ctx.prec,flush=True)
    cases = [
        ('beta_3_over_7',Q(3,7),Q(3,7),Q(1,20000)),
        ('beta_upper',BU,BU,Q(1,10**13)),
    ]
    seed_records=[]
    for name,bl,bu,target in cases:
        lo,hi,seeds=coefficients(bl,bu)
        p,e=lower_polynomial(lo,hi)
        assert e < Q(1,10**20)
        certify_positive(p,Q(0),Q(4),target,name)
        assert bl*bl+2*bl-1>0  # positive beta-mixture tail for z>=2
        seed_records.append({'name':name,'beta':[str(bl),str(bu)],'seeds':seeds,
                             'tail_bound':str(e)})

    lo,hi,seeds=coefficients(BL,BL)
    # Upper polynomial, with an absolute uniform remainder bound.
    p=[((-1)**n)*(hi[n] if n%2==0 else lo[n]) for n in range(N)]
    p[0]+=tail_bound(hi,0)
    x=Q(511574841194,10**12)
    upper=eval_poly(p,x)
    assert upper < -Q(1,10**13)
    print('beta_lower: g_beta(x) < -1/10^13 at x =',x,flush=True)
    seed_records.append({'name':'beta_lower','beta':[str(BL),str(BL)],'seeds':seeds,
                         'negative_x':str(x),'upper_bound':str(upper)})

    # Uniform parameter-box certificates isolate the only possible contact.
    lo,hi,seeds=coefficients(BL,BU)
    p,e=lower_polynomial(lo,hi)
    certify_positive(p,Q(0),Q(1,2),Q(1,10**6),'parameter_box_left')
    certify_positive(p,Q(53,100),Q(4),Q(1,10**6),'parameter_box_right')
    p2,e2=lower_polynomial(lo,hi,2)
    assert e2 < Q(1,10**16)
    certify_positive(p2,Q(1,2),Q(53,100),Q(1,20),'parameter_box_second_derivative')
    seed_records.append({'name':'parameter_box','beta':[str(BL),str(BU)],'seeds':seeds,
                         'tail_bound':str(e),'second_derivative_tail':str(e2)})

    assert 21**5 < 2**22
    assert 21**7 < 2**31
    assert BU < Q(3,7) < Q(13,30)
    print('ALL CERTIFICATES PASSED',flush=True)
    result={'precision_bits':ctx.prec,'N':N,'gamma_decimal_digits':DIGITS,
            'parameters':seed_records,'certificates':records}
    out=Path(__file__).with_suffix('.json')
    out.write_text(json.dumps(result,indent=2)+'\n')
    print('Certificate record:',out,flush=True)


if __name__=='__main__':
    main()
\end{filecontents*}
\begin{filecontents*}[overwrite]{leroy_independent_certificate.py}
"""Independent compact positivity certificate for beta=3/7.

The certified raw density is
 g(x) = sum_{n>=0} (-x)^n/(n! Gamma(3/7-(n+1)/3)^2).
The normalized inverse Mellin density is Gamma(3/7)^2*g(x).

Requires python-flint. All interval inputs and comparison thresholds are
exact rational numbers. The proof of the tail estimate and the interval
coverage argument is in the accompanying manuscript.
"""
from flint import arb, arb_poly, fmpq, ctx
import flint


def main():
    ctx.prec = 256
    beta = fmpq(3, 7)
    N = 90
    # Three initial Gamma values, followed by exact rational recurrences.
    coeff = [arb(0) for _ in range(N + 3)]
    coeff[0] = arb(beta - fmpq(1, 3)).rgamma() ** 2
    coeff[1] = -(arb(beta - fmpq(2, 3)).rgamma() ** 2)
    coeff[2] = arb(beta - 1).rgamma() ** 2 / 2
    for n in range(N):
        coeff[n + 3] = -coeff[n] * arb(
            (n + 4 - 3 * beta) ** 2 / (9 * (n + 1) * (n + 2) * (n + 3))
        )
    P = arb_poly(coeff[:N])
    R = fmpq(2, 25)
    # At x <= 4 and n >= N, the positive absolute term ratio is <= R.
    assert fmpq(64 * (N + 3), 9 * (N + 1) * (N + 2)) < R
    global_target = fmpq(1, 20000)
    records = [
        (0, 2048, fmpq(9, 100000), fmpq(1, 7)),
        (1, 64, fmpq(9, 1000), fmpq(1, 2)),
        (2, 64, fmpq(3, 50), fmpq(3)),
        (3, 512, fmpq(1, 20), fmpq(29)),
    ]
    print('python-flint:', flint.__version__)
    print('precision:', ctx.prec, 'bits; truncation:', N, 'terms')
    for left, pieces, grid_lower, derivative_upper in records:
        X = left + 1
        seeds = [abs(coeff[n]) * X ** n for n in range(N, N + 3)]
        tail = sum(seeds, arb(0)) / arb(1 - R)
        # For n=N,N+1,N+2 and k>=0, (n+3k)t_(n+3k)
        # <= (n+3k)t_n R^k. Summing gives the derivative tail below.
        derivative_tail = sum((
            seeds[j] * arb(fmpq(N + j, 1) / (1 - R) + 3 * R / (1 - R) ** 2) / X
            for j in range(3)
        ), arb(0))
        derivative_bound = sum((
            n * abs(coeff[n]) * X ** (n - 1) for n in range(1, N)
        ), arb(0)) + derivative_tail
        assert derivative_bound < arb(derivative_upper)
        minimum = None
        minimizer = None
        for j in range(pieces + 1):
            x = fmpq(left, 1) + fmpq(j, pieces)
            lower = P(arb(x)) - tail
            assert lower > arb(grid_lower), (left, x, lower)
            if minimum is None or float(lower.mid()) < float(minimum.mid()):
                minimum, minimizer = lower, x
        # Every x in [left,left+1] lies within 1/(2*pieces) of a grid point.
        rigorous_margin = grid_lower - derivative_upper / (2 * pieces)
        assert rigorous_margin > global_target
        print('interval:', (left, X), '; grid points:', pieces + 1)
        print('  derivative majorant:', derivative_bound)
        print('  asserted derivative upper bound:', derivative_upper)
        print('  uniform absolute tail:', tail)
        print('  smallest midpoint-selected grid lower ball:', minimum, 'at', minimizer)
        print('  every grid value strictly greater than:', grid_lower)
        print('  certified whole-interval lower bound:', rigorous_margin)
    print('CERTIFIED: g(x) > 1/20000 for every real x in [0,4].')
    print('Hence q_(3/7)(x) > Gamma(3/7)^2/20000 on [0,4].')


if __name__ == '__main__':
    main()
\end{filecontents*}

\begin{document}
\maketitle
\begin{abstract}
For the Le Roy function
\(F_{\alpha,\beta}^{(\gamma)}(z)=\sum_{n\ge0}z^n/\Gamma(\beta+\alpha n)^\gamma\),
we give a completely specified counterexample to Simon's proposed
equivalence between complete monotonicity and logarithmic infinite
divisibility. The function \(F_{1/3,3/7}^{(2)}(-x)\) is completely
monotone, whereas its unique representing probability law is not
logarithmically infinitely divisible. A beta-product integral proves
positivity of the density for \(x^3\ge54\); an exact rational
Bernstein certificate, with three rigorously enclosed gamma constants,
proves positivity on \([0,4]\). The failure of the proposed necessary
condition reduces to \(21^7<2^{31}\).
We also enclose the complete-monotonicity threshold at
\((\alpha,\gamma)=(1/3,2)\) in the interval
\((0.428292245841,0.428292245842)\), prove that its density has
exactly one zero, of order two, obtain an explicit interval of
counterexamples, and prove a \(3/2\)-power law for the negative mass
born below the boundary. All certificates are included in the single source
file. The qualitative refutation was already present in corpus item
[116]; the contributions here are an explicit independent proof,
a rigorous threshold enclosure, and uniqueness of the critical contact.
\end{abstract}

\section{The external conjecture and the precise result}

A smooth function \(f:(0,\infty)\to\R\) is completely monotone
(CM) when \((-1)^n f^{(n)}(x)\ge0\) for every integer \(n\ge0\)
and every \(x>0\). A positive random variable is logarithmically
infinitely divisible (log-ID) when its logarithm has an infinitely
divisible distribution on \(\R\).

Simon \cite[Section 3]{Simon} proposed the following equivalence for
all positive \(\alpha,\beta,\gamma\):
\begin{equation}\label{eq:external}
 F_{\alpha,\beta}^{(\gamma)}(-x)\text{ is CM}
 \quad\Longleftrightarrow\quad
 z+\gamma(z^{\beta-\alpha}-z^\beta)\le1
 \quad(0<z<1).
\end{equation}
The scalar condition gives a nonnegative Levy kernel and is sufficient.
Its claimed necessity is the part refuted here. The same paper explains
the equivalent formulation in terms of log-ID of the representing law.

Throughout this manuscript, \(\alpha=1/3\) and \(\gamma=2\). Put
\begin{equation}\label{eq:global-series}
 g_\beta(x)=\sum_{n=0}^{\infty}
 \frac{(-x)^n}{n!\,\Gamma(\beta-(n+1)/3)^2},\qquad
 q_\beta(x)=\Gamma(\beta)^2 g_\beta(x).
\end{equation}
We prove below that \(q_\beta\) is an integrable signed density of
mass one, with the desired Mellin transform. The series is entire in
\(x\); no positivity is assumed in its definition.

\begin{theorem}[A small rational counterexample]\label{thm:explicit}
There exists a positive random variable \(X\), with a strictly positive
density on \((0,\infty)\), such that
\[
 \E X^s=\frac{\Gamma(1+s)\Gamma(3/7)^2}
 {\Gamma(3/7+s/3)^2}\quad(\Rea s>-1),
 \qquad
 \E e^{wX}=\Gamma(3/7)^2F_{1/3,3/7}^{(2)}(w)
 \quad(w\in\C).
\]
Consequently \(F_{1/3,3/7}^{(2)}(-x)\) is CM. However, \(\log X\)
is not infinitely divisible, and the right-hand side of
\eqref{eq:external} fails at \(z=1/64\).
\end{theorem}

Define the exact rational numbers
\begin{equation}\label{eq:box}
 \bL=\frac{428292245841}{10^{12}},\qquad
 \bU=\frac{428292245842}{10^{12}}.
\end{equation}

\begin{theorem}[Boundary enclosure and a unique double contact]
\label{thm:boundary}
There is a unique number \(\beta_*\in(\bL,\bU)\) for which
\[
 F_{1/3,\beta}^{(2)}(-x)\text{ is CM}
 \quad\Longleftrightarrow\quad\beta\ge\beta_*,
 \qquad \beta>0.
\]
The critical density \(q_{\beta_*}\) has exactly one zero on
\([0,\infty)\). This zero \(x_*\) belongs to \((1/2,53/100)\),
and
\[
 q_{\beta_*}(x_*)=q'_{\beta_*}(x_*)=0,
 \qquad q''_{\beta_*}(x_*)>\frac{\Gamma(\beta_*)^2}{20}.
\]
In particular, its order is exactly two. The equations
\(g_\beta(x)=\partial_xg_\beta(x)=0\) have exactly one solution
in \([\bL,\bU]\times[1/2,53/100]\).
For every \(\beta\in[\bU,13/30]\), the function is CM and its
representing law is not log-ID.
\end{theorem}

\paragraph{What is new relative to the supplied papers.}
Corpus item [116] \cite{Corpus116} already proves a qualitative
strict separation for all \(\gamma>1\), \(0<\alpha<1/\gamma\),
without specifying a rational counterexample. The preceding equivalence
paper \cite{Bridges} gives, among other results, the lower bound
\(\beta_*>\bL\). This manuscript does not count the same external
conjecture twice. Its additional content is a short explicit rational
triple, an independent beta-product proof of its global positivity,
the matching upper bound \(\beta_*<\bU\), and the uniqueness and
quadratic nature of the boundary contact. Neither the qualitative
strict-gap theorem nor unproved uniform saddle estimates are used in
the proof below.

\paragraph{Proof architecture.}
An exact beta-product identity handles the infinite interval. On the
remaining compact interval, the entire density series has a rational
three-step recurrence after three gamma constants are enclosed. We
bound the infinite tail and convert a finite lower polynomial into
Bernstein form. All coefficient sign decisions are exact rational
comparisons. This supplies the positive probability density; failure
of log-ID then follows from a separate integer inequality.

% Analytical fragment only.  Labels use the prefix mix:.
% Required bibliography keys: Simon, DLMF.
% The finite positivity certificate is supplied by the computational section.

\section{An explicit signed beta mixture and its positivity}
\label{mix:section}

Throughout this section, \(\beta>0\), and we write
\[
 F_\beta(w)=F^{(2)}_{1/3,\beta}(w)
       =\sum_{n=0}^{\infty}\frac{w^n}{\Gamma(\beta+n/3)^2},
 \qquad
 \mathcal F_\beta(w)=\Gamma(\beta)^2F_\beta(w).
\]
The normalization is \(\mathcal F_\beta(0)=1\).  The construction
below is valid before positivity has been established: its output
is initially a finite signed density.  In particular, no existence
of a random variable having the desired Mellin transform is assumed.

\begin{lemma}[A signed density from two beta variables]
\label{mix:signed-density}
Let \(B_1,B_2,G\) be independent, with
\[
 B_1\sim\operatorname{Beta}\!\left(\frac13,\beta+\frac23\right),
 \qquad
 B_2\sim\operatorname{Beta}\!\left(\frac23,\beta+\frac13\right),
 \qquad G\sim\operatorname{Exp}(1),
\]
and put \(B=B_1B_2\).  Define
\[
 P_\beta(y)=y^2+(2\beta-3)y+(\beta-1)^2
\]
and, for \(z>0\),
\begin{equation}\label{mix:Z-density}
 h_\beta(z)=\frac1{\beta^2}\,
   \mathbb E\!\left[B^{-1}e^{-z/B}P_\beta(z/B)\right].
\end{equation}
This is a smooth finite signed density of total mass one.  Its Mellin
integral is absolutely convergent for \(\Re s>-1/3\), and equals
\begin{align}
 \int_0^\infty z^s h_\beta(z)\,dz
 &=\frac{\Gamma(\beta)^2
       \Gamma(s+\frac13)\Gamma(s+\frac23)\Gamma(s+1)}
       {\Gamma(\frac13)\Gamma(\frac23)\Gamma(s+\beta)^2}
 \label{mix:Z-mellin}\\
 &=27^{-s}\frac{\Gamma(1+3s)\Gamma(\beta)^2}
                          {\Gamma(\beta+s)^2}.
 \nonumber
\end{align}
Consequently the change of variables \(z=x^3/27\) gives a smooth
signed density on \((0,\infty)\),
\begin{equation}\label{mix:X-density}
 q_\beta(x)=\frac{x^2}{9}\,h_\beta(x^3/27),
\end{equation}
whose absolutely convergent Mellin transform is
\begin{equation}\label{mix:X-mellin}
 M_\beta(s):=\int_0^\infty x^s q_\beta(x)\,dx
       =\frac{\Gamma(1+s)\Gamma(\beta)^2}
                    {\Gamma(\beta+s/3)^2},
       \qquad \Re s>-1.
\end{equation}
The reciprocals of Gamma functions in these formulas are understood
as entire functions, so zeros of a denominator cause no ambiguity.
\end{lemma}

\begin{proof}
For \(\sigma>-1/3\), the beta integrals give
\begin{equation}\label{mix:beta-product-mellin}
 \mathbb E B^s
 =\frac{\Gamma(s+\frac13)\Gamma(s+\frac23)
             \Gamma(\beta+1)^2}
       {\Gamma(\frac13)\Gamma(\frac23)
             \Gamma(s+\beta+1)^2},
       \qquad \Re s=\sigma.
\end{equation}
In particular, \(\mathbb E B^\sigma<\infty\).  Substituting
\(z=By\) and applying Tonelli to absolute values yields
\[
 \int_0^\infty z^\sigma|h_\beta(z)|\,dz
 \le\frac{\mathbb E B^\sigma}{\beta^2}
       \int_0^\infty y^\sigma e^{-y}|P_\beta(y)|\,dy<\infty.
\]
This justifies Fubini also for complex \(s\) in the stated half-plane.
The elementary Gamma integral is
\begin{align*}
 \int_0^\infty y^s e^{-y}P_\beta(y)\,dy
 &=\Gamma(s+1)\bigl((s+1)(s+2)
                +(2\beta-3)(s+1)+(\beta-1)^2\bigr)\\
 &=\Gamma(s+1)(s+\beta)^2.
\end{align*}
Combining this with \eqref{mix:beta-product-mellin} and
\(\Gamma(\beta+1)=\beta\Gamma(\beta)\) proves the first equality
in \eqref{mix:Z-mellin}.  Gauss's multiplication formula with
multiplier three, \cite[\S5.5(iii), Eq.~5.5.6]{DLMF}, gives the
second equality.  At \(s=0\) the integral is one.  Smoothness on
every compact subinterval of \((0,\infty)\) follows by differentiating
under the expectation: the exponential suppresses every negative
power of \(B\) when \(z\) stays bounded away from zero.

Finally, \(dz/dx=x^2/9\), and
\[
 \int_0^\infty x^s q_\beta(x)\,dx
 =27^{s/3}\int_0^\infty z^{s/3}h_\beta(z)\,dz.
\]
The factors \(27^{s/3}\) and \(27^{-s/3}\) cancel exactly, giving
\eqref{mix:X-mellin} and its absolute convergence for \(\Re s>-1\).
\end{proof}

The signed-density formula can also be expressed by a differential
operator acting on an ordinary positive density.  Indeed, \(W=BG\)
has density
\[
 f_W(z)=\mathbb E[B^{-1}e^{-z/B}].
\]
For \(L_\beta=-z\,d/dz+\beta-1\), differentiation under this
expectation gives
\begin{equation}\label{mix:differential-mixture}
              h_\beta=\beta^{-2}L_\beta^2f_W.
\end{equation}
The preceding proof establishes the Mellin identity directly, so
\eqref{mix:differential-mixture} does not require any unproved
integration-by-parts boundary conditions.

\begin{lemma}[Entire transform and the positivity criterion]
\label{mix:entire-transform}
For every \(a>0\),
\[
              \int_0^\infty e^{ax}|q_\beta(x)|\,dx<\infty.
\]
Consequently, for every \(w\in\mathbb C\),
\begin{equation}\label{mix:mgf}
 \int_0^\infty e^{wx}q_\beta(x)\,dx
       =\Gamma(\beta)^2\sum_{n=0}^{\infty}
              \frac{w^n}{\Gamma(\beta+n/3)^2}
       =\mathcal F_\beta(w).
\end{equation}
In particular,
\begin{equation}\label{mix:cm-criterion}
 F_\beta(-\lambda)\text{ is completely monotone for }\lambda>0
 \quad\Longleftrightarrow\quad
 q_\beta(x)\ge0\text{ for every }x>0.
\end{equation}
\end{lemma}

\begin{proof}
For \(y\ge0\), the function
\(y|P_\beta(y)|e^{-y/2}\) is bounded.  Since \(y=z/B\ge z\),
\eqref{mix:Z-density} therefore gives, for \(z\ge1\),
\[
        |h_\beta(z)|\le C_\beta z^{-1}e^{-z/2}.
\]
Thus \(|q_\beta(x)|\le C'_\beta x^{-1}e^{-x^3/54}\) for
\(x\ge3\).  Near zero, absolute integrability already follows
from Lemma~\ref{mix:signed-density}.  This proves all the asserted
absolute exponential moments.  Expanding the exponential and using
\eqref{mix:X-mellin} at nonnegative integers now proves
\eqref{mix:mgf} by absolute convergence.

If \(q_\beta\ge0\), its mass is one and
\eqref{mix:mgf} is a probability Laplace transform.  Conversely,
complete monotonicity and \(\mathcal F_\beta(0)=1\) give a probability
measure on \([0,\infty)\) by Bernstein's theorem.  Its Laplace
transform equals that of the finite signed measure
\(q_\beta(x)\,dx\).  Uniqueness of finite-measure Laplace transforms
identifies the two measures.  Continuity of \(q_\beta\) then gives
pointwise nonnegativity.  This also proves uniqueness of the
representing probability law without invoking a moment-indeterminacy
argument.
\end{proof}

\begin{lemma}[The density series, justified by a finite contour shift]
\label{mix:series-ode}
For every \(\beta>0\), the signed density extends to an entire
function of \(x\), with the convergent series
\begin{equation}\label{mix:density-series}
 q_\beta(x)=\Gamma(\beta)^2 g_\beta(x),
 \qquad
 g_\beta(x)=\sum_{n=0}^{\infty}
      \frac{(-x)^n}{n!\,\Gamma(\beta-(n+1)/3)^2}.
\end{equation}
It satisfies the regular differential equation
\begin{equation}\label{mix:ode}
 9q_\beta'''(x)+x^2q_\beta''(x)
        +(9-6\beta)xq_\beta'(x)+(4-3\beta)^2q_\beta(x)=0.
\end{equation}
\end{lemma}

\begin{proof}
Stirling's estimate on vertical strips, \cite[\S5.11]{DLMF}, gives
\[
 M_\beta(c+it)
 =O\!\left((1+|t|)^{c/3+3/2-2\beta}e^{-\pi|t|/6}\right),
 \qquad |t|\longrightarrow\infty,
\]
uniformly for \(c\) in any fixed bounded real interval.  Mellin
inversion consequently holds absolutely on the line \(c=0\):
\[
 q_\beta(x)=\frac1{2\pi i}\int_{-i\infty}^{i\infty}
                  M_\beta(s)x^{-s-1}\,ds.
\]
Move this contour only as far as \(\Re s=-9/2\).  The horizontal
integrals tend to zero by the displayed estimate.  The only possible
poles crossed are \(-1,-2,-3,-4\), and their residues give
\[
 q_\beta(x)=\Gamma(\beta)^2
    \sum_{n=0}^{3}\frac{(-x)^n}{n!\,\Gamma(\beta-(n+1)/3)^2}
       +R_\beta(x).
\]
A reciprocal-Gamma zero can cancel a pole, in which case the same
residue formula simply gives zero.  Differentiating the remaining
absolutely convergent integral shows
\[
 R_\beta^{(k)}(x)=O(x^{7/2-k}),
           \qquad 0\le k\le3,\quad x\downarrow0.
\]
Hence \(q_\beta\) has a \(C^3\) extension to zero and, for
\(k=0,1,2\),
\begin{equation}\label{mix:initial-values}
 q_\beta^{(k)}(0)=
       \frac{\Gamma(\beta)^2(-1)^k}
                    {\Gamma(\beta-(k+1)/3)^2}.
\end{equation}

Gamma's recurrence gives the meromorphic identity
\[
 (s+3\beta-3)^2M_\beta(s)
       =9s(s-1)(s-2)M_\beta(s-3).
\]
For \(\Re s>2\), the Mellin transform of the left side of
\eqref{mix:ode}, evaluated by integration by parts, is therefore
\[
 -9s(s-1)(s-2)M_\beta(s-3)
            +(s+3\beta-3)^2M_\beta(s)=0.
\]
All the boundary terms vanish.  At zero this follows from the
\(C^3\) extension, and at infinity it follows either from the beta
mixture or from differentiated Mellin inversion on arbitrarily large
positive abscissae.  Mellin-transform uniqueness proves
\eqref{mix:ode} on \((0,\infty)\), and continuity extends it to zero.

The coefficients
\[
 a_n=\frac{\Gamma(\beta)^2(-1)^n}
                  {n!\,\Gamma(\beta-(n+1)/3)^2}
\]
satisfy, including when a coefficient vanishes,
\begin{equation}\label{mix:coefficient-recurrence}
 a_{n+3}=-\frac{(n+4-3\beta)^2}
                    {9(n+1)(n+2)(n+3)}\,a_n.
\end{equation}
Each of the three residue-class series is entire by the ratio test.
Their sum satisfies \eqref{mix:ode} and the three initial conditions
\eqref{mix:initial-values}.  The leading coefficient of
\eqref{mix:ode} is the nonzero constant nine.  Uniqueness for this
ordinary initial-value problem identifies its entire series solution
with \(q_\beta\) for all \(x\ge0\).  Thus the proof uses one finite
contour shift and does not close a Mellin contour through infinitely
many poles.
\end{proof}

\begin{corollary}[Every nonnegative contact is quadratic]
\label{mix:quadratic-contact}
If \(q_\beta\ge0\) on \((0,\infty)\) and
\(q_\beta(x_0)=0\) at \(x_0>0\), then
\[
                q_\beta'(x_0)=0,
                \qquad q_\beta''(x_0)>0.
\]
\end{corollary}

\begin{proof}
Nonnegativity gives \(q_\beta'(x_0)=0\) and
\(q_\beta''(x_0)\ge0\).  If the second derivative were zero,
uniqueness for the regular third-order equation \eqref{mix:ode}
would give \(q_\beta\equiv0\), contrary to its total mass one.
\end{proof}

\begin{lemma}[An elementary positive tail]
\label{mix:positive-tail}
For every \(\beta\ge\sqrt2-1\),
\[
                   q_\beta(x)>0
          \qquad\text{whenever }x\ge\sqrt[3]{54}.
\]
In particular this holds for every \(\beta\ge5/12\).  At the
principal rational parameter \(\beta=3/7\),
\[
 P_{3/7}(y)=y^2-\frac{15}{7}y+\frac{16}{49},
 \qquad P_{3/7}(2)=\frac2{49}>0.
\]
\end{lemma}

\begin{proof}
For \(y\ge2\),
\[
 P_\beta(2)=\beta^2+2\beta-1\ge0,
 \qquad P_\beta'(y)=2y+2\beta-3\ge1+2\beta>0.
\]
Since \(0<B<1\) almost surely, \(z\ge2\) implies \(z/B>2\).
The integrand in \eqref{mix:Z-density} is then strictly positive.
Equation~\eqref{mix:X-density} completes the proof.  The rational
sufficient bound follows also from
\(P_{5/12}(2)=1/144>0\).
\end{proof}

\begin{lemma}[Transfer to all larger beta parameters]
\label{mix:beta-transfer}
Suppose \(q_\beta\ge0\), and let \(\beta'>\beta\).  Take independent
\[
 U_1,U_2\sim\operatorname{Beta}(\beta,\beta'-\beta),
 \qquad R=(U_1U_2)^{1/3},
\]
independent also of a random variable \(X_\beta\) with density
\(q_\beta\).  Then
\begin{equation}\label{mix:transfer-law}
                    X_{\beta'}\overset d=X_\beta R
\end{equation}
has density \(q_{\beta'}\).  In particular \(F_{\beta'}(-\lambda)\)
is completely monotone.  If \(q_\beta\) is strictly positive on a
terminal interval of the positive half-line, then \(q_{\beta'}\)
is strictly positive at every positive point.
\end{lemma}

\begin{proof}
The ordinary beta integral gives
\[
 \mathbb E R^s=
    \left(\frac{\Gamma(\beta')\Gamma(\beta+s/3)}
                {\Gamma(\beta)\Gamma(\beta'+s/3)}\right)^2,
                \qquad \Re s>-3\beta.
\]
Multiplication by \(M_\beta(s)\) yields \(M_{\beta'}(s)\).
Taking \(s=it\), uniqueness of Fourier transforms after the change
of variables \(x=e^y\) identifies the product law with the finite
signed measure \(q_{\beta'}(x)\,dx\); hence the latter is positive.
Equivalently, its continuous density is
\[
                q_{\beta'}(x)
                    =\mathbb E[R^{-1}q_\beta(x/R)].
\]
The law of \(R\) has support \((0,1)\).  For every fixed \(x>0\),
the event that \(x/R\) lies in the positive terminal interval has
positive probability.  The last formula is therefore strictly
positive.  Complete monotonicity follows from
Lemma~\ref{mix:entire-transform}.
\end{proof}

\begin{lemma}[An exact obstruction to logarithmic infinite divisibility]
\label{mix:not-id}
If \(0<\beta\le13/30\) and \(q_\beta\ge0\), a random variable
\(X_\beta\) with density \(q_\beta\) is not logarithmically
infinitely divisible.
\end{lemma}

\begin{proof}
The Gamma logarithm formula used in \cite[Section~2]{Simon} gives
\[
 \log M_\beta(it)=ic_\beta t
       +\int_0^\infty
          (e^{-itr}-1+itr)\frac{\phi_\beta(r)}{r}\,dr,
\]
where \(c_\beta=\psi(1)-\frac23\psi(\beta)\) and
\begin{equation}\label{mix:levy-kernel}
 \phi_\beta(r)=\frac{e^{-r}}{1-e^{-r}}
                   -2\frac{e^{-3\beta r}}{1-e^{-3r}}.
\end{equation}
The integral is absolutely convergent after compensation, and
\(\int_0^\infty r|\phi_\beta(r)|\,dr<\infty\).
If \(\log X_\beta\) were infinitely divisible, its L\'evy measure
would have to be nonnegative, forcing \(\phi_\beta\ge0\).
For completeness, this necessity follows directly from uniqueness:
the negative second derivative of the displayed exponent is the
Fourier transform of the finite signed measure supported at \(-r\)
with mass element \(r\phi_\beta(r)\,dr\).  On the other hand, the
negative second derivative of an infinitely divisible exponent with
finite variance is the Fourier transform of its nonnegative
variance measure, including any Gaussian mass at zero.  Fourier
uniqueness thus forces \(r\phi_\beta(r)\,dr\) to be nonnegative.
The required finite variance of \(\log X_\beta\) follows from
\eqref{mix:X-mellin} at small positive and negative real arguments.

Put \(u=e^{-r}\).  Nonnegativity in
\eqref{mix:levy-kernel} is equivalent to
\begin{equation}\label{mix:scalar-id-condition}
                 2u^{3\beta-1}\le1+u+u^2,
                         \qquad 0<u<1.
\end{equation}
At \(u=1/4\) and \(\beta\le13/30\), the left side is at least
\(2^{2/5}\), whereas the right side is \(21/16\).  The strict
reverse inequality is certified by the exact integer comparison
\[
           2^{22}=4\,194\,304
                      >4\,084\,101=21^5.
\]
Consequently \(\phi_\beta(\log4)<0\), proving the assertion.
At \(\beta=3/7\) one can alternatively use the particularly simple
comparison
\[
          21^7=1\,801\,088\,541
                     <2\,147\,483\,648=2^{31}.
\]
\end{proof}

\begin{corollary}[An explicit interval of counterexamples]
\label{mix:explicit-family}
For every
\[
                     \frac37\le\beta\le\frac{13}{30},
\]
the function \(F^{(2)}_{1/3,\beta}(-\lambda)\) is completely
monotone on \((0,\infty)\), and its normalized representing law
has a strictly positive density and is not logarithmically infinitely
divisible.  Complete monotonicity in fact holds for every
\(\beta\ge3/7\).
\end{corollary}

\begin{proof}
The finite certificate accompanying this proof establishes
\begin{equation}\label{mix:compact-certificate}
                  g_{3/7}(x)>\frac1{20000}
                           \qquad(0\le x\le4).
\end{equation}
It uses the degree-89 truncation of \eqref{mix:density-series},
a rigorous three-step recurrence remainder, and four rational
Bernstein subintervals.  The certificate and its independent
verification are supplied in the computational section; the
analytical identity needed by that certificate is proved in
Lemma~\ref{mix:series-ode}.
Since \(\sqrt[3]{54}<4\), \eqref{mix:compact-certificate} and
Lemma~\ref{mix:positive-tail} cover the whole positive half-line.
Lemmas~\ref{mix:entire-transform} and \ref{mix:beta-transfer} prove
complete monotonicity and strict density positivity for every larger
parameter.  Lemma~\ref{mix:not-id} gives the stated interval on which
logarithmic infinite divisibility fails.
\end{proof}


\section{Finite certificates for an infinite interval problem}
\label{sec:certificate}

This section supplies the compact-interval input to the analytical
construction. It also proves the lower boundary obstruction and the
uniform estimates used to isolate its contact point. The computations
are proofs by interval enclosure and exact polynomial inequalities,
rather than tests of sampled approximate values.

\subsection{Three gamma constants and a rational recurrence}

Write \(g_\beta(x)=\sum_{n\ge0}(-1)^n a_n(\beta)x^n\), where
\[
 a_n(\beta)=\frac1{n!\Gamma(\beta-(n+1)/3)^2}\ge0.
\]
The gamma recurrence gives, exactly,
\begin{equation}\label{eq:coefficient-rec}
 a_{n+3}(\beta)=a_n(\beta)
 \frac{(n+4-3\beta)^2}{9(n+1)(n+2)(n+3)}.
\end{equation}
All coefficients are therefore determined by \(a_0,a_1,a_2\).
For \(\beta\in[\ell,u]\subset(1/3,1/2)\), the positive multiplier
in \eqref{eq:coefficient-rec} decreases as \(\beta\) increases.
If \(A_n^-\le a_n(\beta)\le A_n^+\) throughout that parameter
interval, valid next enclosures are
\[
 A_{n+3}^-=A_n^-\frac{(n+4-3u)^2}{9(n+1)(n+2)(n+3)},\quad
 A_{n+3}^+=A_n^+\frac{(n+4-3\ell)^2}{9(n+1)(n+2)(n+3)}.
\]
The program first encloses the three reciprocal-gamma squares with
384-bit Arb arithmetic \cite{Johansson}, then explicitly verifies
strict rational lower and upper bounds with denominator \(10^{80}\).
All subsequent operations use exact rational numbers.

\subsection{A uniform remainder bound, including derivatives}

Take \(N=90\). For \(d\in\{0,1,2\}\) define
\(T_n^{(d)}=a_n(\beta)n^{\underline d}4^{n-d}\), with
\(n^{\underline d}=n!/(n-d)!\). For \(n\ge90\),
\begin{align}
 \frac{T_{n+3}^{(d)}}{T_n^{(d)}}
 &=\frac{64(n+4-3\beta)^2}
 {9(n+1-d)(n+2-d)(n+3-d)}\notag\\
 &\le\frac{64(n+3)^2}{9(n-1)n(n+1)}
 \le\frac{64\cdot93^2}{9\cdot89\cdot90\cdot91}
 <\frac1{10}.\label{eq:tail-ratio}
\end{align}
The rational function in the middle is decreasing for \(n\ge90\);
this follows, for example, from its logarithmic derivative
\(2/(n+3)-1/(n-1)-1/n-1/(n+1)<0\).
Summing the three residue classes modulo three gives
\begin{equation}\label{eq:tail-cert}
 \sup_{0\le x\le4}\left|
 g_\beta^{(d)}(x)-\sum_{n=d}^{89}(-1)^n
 a_n(\beta)n^{\underline d}x^{n-d}\right|
 \le E_d:=\frac{10}{9}
 \sum_{n=90}^{92}A_n^+n^{\underline d}4^{n-d}.
\end{equation}
Every displayed application satisfies \(E_0<10^{-20}\); the
parameter-box second derivative satisfies \(E_2<10^{-16}\).
These comparisons are themselves checked with exact rational arithmetic.

A rational lower polynomial for \(g_\beta^{(d)}\) is obtained by
using \(A_n^-\) for even \(n\), \(A_n^+\) for odd \(n\), and
subtracting \(E_d\) from the constant coefficient. An upper polynomial
uses the opposite endpoints and adds \(E_d\). Positivity of \(x\)
is the only order property needed for these substitutions.

\subsection{Bernstein positivity and a verifiable cover}

Let \(P\) be one of the rational lower polynomials, of degree at most
\(m\), and let \([a,b]\) be a rational interval. After substitution,
\[
 P(a+(b-a)t)=\sum_{k=0}^{m}v_kt^k
 =\sum_{j=0}^{m}B_j\binom mj t^j(1-t)^{m-j},
 \quad
 B_j=\sum_{k=0}^{j}v_k\frac{\binom jk}{\binom mk}.
\]
The Bernstein basis functions are nonnegative and sum to one on
\([0,1]\). Hence \(B_j>c\) for every \(j\) proves \(P>c\)
throughout \([a,b]\). If this test fails, de Casteljau subdivision
at the midpoint gives the exact Bernstein coefficients on the two
halves. The verifier checks that the final closed intervals cover the
whole requested interval without gaps.

\begin{proposition}[Certified compact bounds]\label{prop:certified}
The following strict inequalities hold:
\begin{align}
 g_{3/7}(x)&>\frac1{20000} &&(0\le x\le4),\label{eq:cert-main}\\
 g_{\bU}(x)&>10^{-13} &&(0\le x\le4),\label{eq:cert-upper}\\
 g_{\bL}\!\left(\frac{511574841194}{10^{12}}\right)
 &<-10^{-13},\label{eq:cert-lower}\\
 g_\beta(x)&>10^{-6}
 &&\left(\beta\in[\bL,\bU],\quad
 x\in[0,1/2]\cup[53/100,4]\right),\label{eq:cert-outside}\\
 g_\beta''(x)&>\frac1{20}
 &&\left(\beta\in[\bL,\bU],\quad
 x\in[1/2,53/100]\right).\label{eq:cert-convex}
\end{align}
\end{proposition}

\begin{proof}
Apply the preceding rational enclosures, \eqref{eq:tail-cert}, and
Bernstein coefficient test. The complete executable verifier appears
in Appendix~\ref{app:code}; the exact inputs are the displayed rational
numbers. Its successful output is summarized in Table~\ref{tab:cert}.
For \eqref{eq:cert-lower}, evaluate the upper polynomial at the stated
rational point. There is no floating-point sign decision in this step.
\end{proof}

\begin{table}[ht]
\centering\small
\begin{tabular}{@{}llll@{}}
\toprule
Parameter or box & Interval & Certified lower bound & Pieces\\
\midrule
\(\beta=3/7\) & \([0,4]\) & \(g>1/20000\) & 4\\
\(\beta=\bU\) & \([0,4]\) & \(g>10^{-13}\) & 18\\
\(\beta\in[\bL,\bU]\) & \([0,1/2]\) & \(g>10^{-6}\) & 1\\
\(\beta\in[\bL,\bU]\) & \([53/100,4]\) & \(g>10^{-6}\) & 1\\
\(\beta\in[\bL,\bU]\) & \([1/2,53/100]\) & \(g''>1/20\) & 1\\
\bottomrule
\end{tabular}
\caption{Exact rational Bernstein certificates after gamma enclosures.
The degree is 89 for \(g\), and 87 for \(g''\).}
\label{tab:cert}
\end{table}

For the small rational counterexample, the four pieces are simply
\([0,1/2]\), \([1/2,1]\), \([1,2]\), and \([2,4]\).
For \(\beta=\bU\), the partition endpoints, in increasing order, are
\begin{gather*}
0,\ \tfrac12,\ \tfrac{65}{128},\ \tfrac{261}{512},\
\tfrac{523}{1024},\ \tfrac{1047}{2048},\ \tfrac{2095}{4096},\
\tfrac{8381}{16384},\ \tfrac{16763}{32768},\ \tfrac{4191}{8192},\\
\tfrac{131}{256},\ \tfrac{33}{64},\ \tfrac{17}{32},\
\tfrac9{16},\ \tfrac58,\ \tfrac34,\ 1,\ 2,\ 4.
\end{gather*}
Thus the near-zero minimum is covered by a finite interval proof; the
argument does not infer global positivity from the decimal value of
a numerical minimizer.

\subsection{An independent compact-interval check}

For \(\beta=3/7\), the same conclusion was checked independently by
a grid with a rigorous derivative bound and a rigorous infinite-tail
bound. On a unit interval subdivided into \(K\) parts, assume every
grid point has \(g>L\) and the whole interval has \(|g'|<M\).
Every real point is within \(1/(2K)\) of a grid point, so the mean
value theorem gives \(g>L-M/(2K)\). A rigorous derivative majorant
on \([0,X]\), for \(X\in\{1,2,3,4\}\), is
\[
 \sum_{n=1}^{89}n a_nX^{n-1}
 +\frac1X\sum_{n=90}^{92}a_nX^n
 \left(\frac{n}{1-R}+\frac{3R}{(1-R)^2}\right),
 \qquad R=\frac2{25}.
\]
Indeed, the undifferentiated three-step ratio is at most
\(64\cdot93/(9\cdot91\cdot92)<R\), and
\(\sum_{k\ge0}(n+3k)R^k=n/(1-R)+3R/(1-R)^2\).
The finite sums and gamma constants are enclosed using Arb. The
verified values are
\begin{center}\small
\begin{tabular}{@{}lllll@{}}
\toprule
Interval & \(K\) & \(L\) & \(M\) & \(L-M/(2K)\)\\
\midrule
\([0,1]\)&2048&\(9/100000\)&\(1/7\)&\(4939/89600000\)\\
\([1,2]\)&64&\(9/1000\)&\(1/2\)&\(163/32000\)\\
\([2,3]\)&64&\(3/50\)&\(3\)&\(117/3200\)\\
\([3,4]\)&512&\(1/20\)&\(29\)&\(111/5120\)\\
\bottomrule
\end{tabular}
\end{center}
All four rational lower bounds exceed \(1/20000\). There are 2692
grid evaluations, including repeated endpoints. This independent
method is included in Appendix~\ref{app:independent}. The Bernstein
certificate already suffices for the theorem.

\section{The boundary and its unique quadratic contact}
\label{sec:boundary}

We now prove Theorem~\ref{thm:boundary} without assuming the
qualitative strict-gap theorem from the corpus. The signed-density
construction and Laplace-transform uniqueness give
\begin{equation}\label{eq:CM-positive}
 F_{1/3,\beta}^{(2)}(-x)\text{ is CM}
 \quad\Longleftrightarrow\quad q_\beta(x)\ge0
 \text{ for every }x>0.
\end{equation}
Indeed, a nonnegative signed density of mass one is a probability
density. Conversely, the Bernstein representing positive measure and
\(q_\beta(x)\dd x\) would have the same Laplace transform and hence
are equal. If \(q_\beta\) is negative at a point, continuity excludes
such a positive measure.

The function \((\beta,x)\mapsto q_\beta(x)\) is jointly continuous
on compact subsets of \((0,\infty)\times[0,\infty)\). Indeed, the
three initial reciprocal-gamma coefficients are continuous in
\(\beta\), and the recurrence has ratio \(O(1/n)\) uniformly on
bounded parameter intervals. The three residue-class series thus
converge uniformly on compact sets. This supplies the compact
continuity used below, including at \(x=0\).

The beta-factor transfer proved above makes the CM set an upper set in
\(\beta>0\). It is nonempty by \eqref{eq:cert-upper} and the
positive tail; it does not contain \(\bL\) by
\eqref{eq:cert-lower}. Let \(\beta_*\) be its infimum. It follows
that \(\beta_*\in[\bL,\bU]\), and pointwise continuity of
\(q_\beta\) in \(\beta\) shows \(q_{\beta_*}\ge0\).
Thus \(\beta_*\) itself is in the CM set.

The lower inequality is strict because \(q_{\bL}\) is negative at
the certified point. For the upper inequality, \(q_{\bU}\) has a
positive minimum on \([0,4]\) by \eqref{eq:cert-upper}. Compact
continuity preserves positivity when \(\beta\) decreases slightly.
The tail remains positive uniformly for all these parameters: they
can be kept above \(5/12\), and
\[
 (5/12)^2+2(5/12)-1=1/144>0.
\]
The beta-mixture tail criterion therefore applies on \(x^3\ge54\),
which overlaps \([0,4]\). A parameter strictly below \(\bU\)
is CM, proving \(\beta_*<\bU\).

If \(q_{\beta_*}\) were positive throughout \([0,4]\), the same
compact-continuity and positive-tail argument would yield a CM
parameter below \(\beta_*\), a contradiction. Hence it has a zero.
The uniform bound \eqref{eq:cert-outside} and the positive tail place
every zero in \((1/2,53/100)\). On that interval,
\eqref{eq:cert-convex} makes \(g_{\beta_*}\) strictly convex. A
nonnegative strictly convex function has at most one zero. At the
zero, its derivative vanishes and its second derivative exceeds
\(1/20\), proving the asserted order and inequality for \(q''\).

To prove uniqueness of the pair solving \(g=\partial_xg=0\) in
the displayed rectangle, suppose such a pair \((\beta,x)\) exists.
Strict convexity makes \(x\) the unique minimum on the contact
interval; the outside and tail bounds then give \(q_\beta\ge0\)
everywhere. Thus \(\beta\ge\beta_*\). If \(\beta>\beta_*\),
beta-factor transfer from \(\beta_*\) makes the new density strictly
positive: the scaling factor has a density positive on \((0,1)\),
and the original density vanishes at only one point. This contradicts
\(g_\beta(x)=0\). Consequently \(\beta=\beta_*\) and
\(x=x_*\).

Finally, \(\beta\ge\bU\) implies CM by beta-factor transfer,
whereas \(\beta\le13/30\) violates the scalar log-ID condition
at \(z=1/64\). This proves the explicit interval of counterexamples.

% Analytical corollary of the certified unique critical contact.
% Labels use contact:.  Depends on mix:X-mellin and mix:density-series.

\subsection{The birth of the unique negative interval}
\label{contact:section}

\begin{theorem}[A transverse quadratic contact]
\label{contact:law}
Let \((\beta_*,x_*)\) be the unique critical pair certified above,
and put
\[
 a=\frac12\partial_x^2q_{\beta_*}(x_*)>0,
 \qquad b=\left.\partial_\beta q_\beta(x_*)\right|_{\beta=\beta_*}.
\]
Then \(b>0\).  For every sufficiently small \(\varepsilon>0\),
the density \(q_{\beta_*-\varepsilon}\) is negative on exactly one
interval \((x_-(\varepsilon),x_+(\varepsilon))\), where
\begin{equation}\label{contact:roots}
 x_\pm(\varepsilon)
      =x_*\pm\sqrt{b/a}\,\varepsilon^{1/2}+O(\varepsilon).
\end{equation}
Its total negative mass satisfies
\begin{equation}\label{contact:negative-mass}
 \int_0^\infty\max\{-q_{\beta_*-\varepsilon}(x),0\}\,dx
       =\frac{4b^{3/2}}{3\sqrt a}\,\varepsilon^{3/2}
                                  +o(\varepsilon^{3/2}).
\end{equation}
\end{theorem}

\begin{proof}
Differentiating \eqref{mix:X-mellin}, using the digamma integral,
and applying Mellin inversion gives
\begin{equation}\label{contact:parameter-generator}
 \partial_\beta q_\beta(x)
  =2\int_0^\infty\frac{e^{-\beta t}}{1-e^{-t}}
      \bigl[e^{t/3}q_\beta(xe^{t/3})-q_\beta(x)\bigr]\,dt.
\end{equation}
Here \(\beta\) may range over a compact neighborhood of
\(\beta_*\).  To justify the calculation, on \(s=i\tau\) the
identity being inverted is
\[
 \partial_\beta M_\beta(s)
  =2M_\beta(s)\int_0^\infty
       \frac{e^{-\beta t}(e^{-st/3}-1)}{1-e^{-t}}\,dt.
\]
The bound \(|e^{-i\tau t/3}-1|\le\min\{2,|\tau|t/3\}\),
together with the uniform exponential decay of \(M_\beta(i\tau)\),
justifies Fubini and differentiation under inversion.  The resulting
real-space integral is also absolutely convergent: its bracket is
\(O(t)\) at zero, while the beta-mixture tail estimate and
\(e^{-\beta t}\) control infinity.  At the critical zero, the
subtracted term vanishes and \(q_{\beta_*}(xe^{t/3})>0\) for
\(x=x_*\) and every \(t>0\).  Thus \eqref{contact:parameter-generator}
proves \(b>0\).

The density is jointly real analytic in \((\beta,x)\) near
\((\beta_*,x_*)\), as follows from \eqref{mix:density-series} and
its locally uniformly convergent coefficient recurrence.  Hence
\[
 q_{\beta_*-\varepsilon}(x_*+h)
    =ah^2-b\varepsilon
         +O\bigl(|h|^3+\varepsilon|h|+\varepsilon^2\bigr).
\]
Set \(\delta=\sqrt\varepsilon\) and \(h=\delta y\).
The quotient
\(\delta^{-2}q_{\beta_*-\delta^2}(x_*+\delta y)\)
extends analytically to \(\delta=0\), with value \(ay^2-b\).
The implicit function theorem at its two simple zeros
\(y=\pm\sqrt{b/a}\) proves \eqref{contact:roots}.
The certified positivity outside \((1/2,53/100)\), together with
the certified strict convexity inside this interval for all parameters
in the enclosing beta box, shows that these are the only two zeros
and that their intervening interval is the whole negative set.
Finally, the same rescaling and the Taylor estimate give
\[
 \varepsilon^{-3/2}
    \int_0^\infty\max\{-q_{\beta_*-\varepsilon}(x),0\}\,dx
       \longrightarrow
       \int_{-\sqrt{b/a}}^{\sqrt{b/a}}(b-ay^2)\,dy
       =\frac{4b^{3/2}}{3\sqrt a},
\]
which proves \eqref{contact:negative-mass}.
\end{proof}

\section{Scope of the conclusion}

The literal universal equivalence in \eqref{eq:external} is false;
one rational triple already suffices to establish that statement.
The interval of examples is a strengthening of the same resolution,
not a list of separate conjectures. The full positive-parameter
classification of \(F_{\alpha,\beta}^{(\gamma)}\) is not supplied by
this result. At \((\alpha,\gamma)=(1/3,2)\), however, the boundary
is uniquely specified by the double-contact equations in a rigorously
isolated rectangle, with a parameter enclosure of width \(10^{-12}\).

The method replaces an unbounded positivity check by two complementary
finite objects: a quadratic polynomial that controls the whole tail,
and a rational Bernstein certificate on a compact interval. Its
mathematical dependency on special-function numerics is restricted
to rigorous enclosures of three gamma constants for each fixed
parameter or parameter interval. The supplied source reproduces the
certificate, including the uniform boundary estimates. No external
peer-review or formal proof-assistant validation is asserted.


\begin{thebibliography}{9}\raggedright

\bibitem{Simon}
T. Simon, \emph{Remark on a Mittag--Leffler function of Le Roy type},
Integral Transforms and Special Functions \textbf{33}(2) (2022),
108--114. \href{https://doi.org/10.1080/10652469.2021.1913138}
{doi:10.1080/10652469.2021.1913138}.
Author version: \url{https://arxiv.org/abs/2012.11018}, v2,
25 March 2021; the conjecture is in Section 3.

\bibitem{DLMF}
NIST Digital Library of Mathematical Functions, Chapter 5,
\S5.5(iii), Eq.~5.5.6 (multiplication formula), and \S5.11
(asymptotic expansions of the gamma function).
\url{https://dlmf.nist.gov/5.5}; \url{https://dlmf.nist.gov/5.11}.

\bibitem{Johansson}
F. Johansson, \emph{Arb: efficient arbitrary-precision midpoint-radius
interval arithmetic}, IEEE Transactions on Computers \textbf{66}(8)
(2017), 1281--1292.
\href{https://doi.org/10.1109/TC.2017.2690633}
{doi:10.1109/TC.2017.2690633}.
\url{https://arxiv.org/abs/1611.02831}.

\bibitem{Corpus116}
Keio Research Notes, \emph{Complete monotonicity beyond logarithmic
infinite divisibility: a counterexample theorem for Le Roy functions},
3 November 2026. Unrefereed research manuscript [116] in the
Verifiable Mathematical Reasoning Corpus,
\url{https://github.com/QEEDDD/Verifiable-Mathematical-Reasoning-Corpus},
snapshot \texttt{c30d40119134a59eda6fa5cc6331f27c8e67d034}.

\bibitem{Bridges}
Keio Research Notes, \emph{Equivalence Bridges from the 130-Manuscript
Corpus: explicit independence, arithmetic obstructions, and a
positivity boundary}, 8 November 2026. Preceding research synthesis,
34 pages; the Le Roy section supplies the earlier one-sided bound
\(\beta_c(1/3,2)>0.428292245841\).

\end{thebibliography}

\clearpage
\appendix
\section{Complete exact-rational certificate}
\label{app:code}

The single \texttt{.tex} source contains both verification programs
in \texttt{filecontents*} environments. Compiling it writes the
programs alongside the compilation output. They require Python 3
and \texttt{python-flint}; the recorded runs used
\texttt{python-flint 0.9.0}. Run
\begin{lstlisting}[language={}]
python3 leroy_explicit_certificate.py
python3 leroy_independent_certificate.py
\end{lstlisting}
The first program checks every inequality in
Proposition~\ref{prop:certified} and writes a JSON record containing
the exact rational gamma enclosures, tail bounds, and covering
intervals. It uses Arb only to obtain and verify the initial rational
enclosures. All polynomial bounds and subdivisions use exact rational
arithmetic. The successful run and an independently executed copy
produced identical JSON records. The audited primary program has
SHA256
\begin{center}\ttfamily\footnotesize
2af3eca6a5979ef1df0188b32f1425a2a0a0813d9ec407ed80c111cf370e2d60.
\end{center}
The following is the complete primary program.

\lstinputlisting{leroy_explicit_certificate.py}

\subsection*{Recorded successful output}
\begin{lstlisting}[language={}]
python-flint 0.9.0 ; Arb precision 384
beta_3_over_7 : 4 leaves; max depth 3 ; exact lower bound exceeds 1/20000
beta_upper : 18 leaves; max depth 17 ; exact lower bound exceeds 1/10000000000000
beta_lower: g_beta(x) < -1/10^13 at x = 255787420597/500000000000
parameter_box_left : 1 leaves; max depth 0 ; exact lower bound exceeds 1/1000000
parameter_box_right : 1 leaves; max depth 0 ; exact lower bound exceeds 1/1000000
parameter_box_second_derivative : 1 leaves; max depth 0 ; exact lower bound exceeds 1/20
ALL CERTIFICATES PASSED
\end{lstlisting}

\clearpage
\section{Independent derivative-and-grid certificate}
\label{app:independent}

This second program verifies only the compact positivity assertion
for \(\beta=3/7\), using a different method. Every sign comparison
and error estimate is a rigorous Arb interval comparison with exact
rational inputs. Floating-point conversion appears only when
selecting a diagnostic grid point to print; it has no role in a proof
decision. The interval-cover argument and all thresholds are proved
in Section~\ref{sec:certificate}.

\lstinputlisting{leroy_independent_certificate.py}

\end{document}
